\section{Introduction}
\label{sec:introduction}

Legal question-answering systems must not only retrieve text with relevant
content but also correctly identify the applicable conditions, legal
effects, and statutory basis for each conclusion. In many cases, the answer
does not appear directly in a single provision but must instead be derived
by connecting multiple rules. An initial condition may give rise to an
intermediate legal state, which may then trigger another effect. In addition
to forward reasoning, users may ask the system to trace causes backward,
identify intermediate events, determine whether a relation is direct, or
assess what would happen if a link in the reasoning chain were absent. Legal
benchmarks such as CaseHOLD and LegalBench have demonstrated the diversity
and difficulty of legal reasoning tasks, while LegalBench-RAG further
emphasizes the role of retrieval in legal question-answering systems
\cite{zheng2021casehold,guha2023legalbench,alami2024legalbenchrags}.

Retrieval-Augmented Generation (RAG) combines a generative language model
with externally retrieved knowledge sources to improve the factuality and
traceability of answers \cite{lewis2020rag}. However, conventional RAG
architectures typically represent a corpus as relatively independent text
passages and rank them primarily according to their semantic similarity to
the query. This approach may work well for single-hop questions, but it
struggles when the required evidence is distributed across multiple rules
or when the order in which the evidence is connected determines the meaning
of the answer. Multi-hop question-answering tasks, notably HotpotQA, have
shown that correctly retrieving individual passages does not guarantee that
a system will recover the complete reasoning path \cite{yang2018hotpotqa}.

Graph-based RAG methods partially address this limitation by representing
entities, events, or knowledge units as nodes and the relations among them
as edges. GraphRAG uses graph structure to support querying and information
synthesis at broad scope, whereas HippoRAG combines graph-based retrieval
with a multi-hop association mechanism
\cite{edge2024graphrag,gutierrez2024hipporag}. Compared with retrieving
independent passages, graph representations allow the system to preserve
the linkage structure among pieces of evidence. Nevertheless, an arbitrary
path through a graph does not necessarily correspond to a valid reasoning
chain. In the legal domain, the system must additionally distinguish among
rule-supported condition--effect relations, direct relations, relations
mediated by intermediate events, and alternative paths leading to the same
outcome.

Structural Causal Models (SCMs) provide a formal framework for describing
the mechanisms that generate variable values and for analyzing the results
of interventions \cite{pearl2009causality}. Recently, several studies have
investigated combining causal reasoning or counterfactual reasoning with RAG
\cite{wang2025causalrag,khadilkar2025causalcounterfactualrag,qin2026cfrag}.
These research directions demonstrate the potential of incorporating causal
structure into retrieval and answer generation. However, applying causal
concepts to law requires careful scope limitations. The relations used in
this study are condition--effect relations operationalized from legal rules,
not causal relations discovered or estimated from observational data.
Accordingly, the resulting model should be understood as a structured
normative reasoning model for checking the consistency of rule chains and
analyzing interventions within the given graph.

On this basis, this study presents a CausalRAG pipeline for question
answering in Vietnamese criminal law. The pipeline begins by normalizing
legal rules into conditions, effects, legal subjects, and statutory bases.
Conditions and effects are mapped to events in a heterogeneous legal causal
graph, in which rules serve as mechanisms connecting condition events to
effect events. In parallel with the graph, the system constructs a vector
memory for rules and events to retrieve query-relevant seed events. From
these seed events, the retriever expands and ranks causal paths of up to two
hops before passing the candidate paths to the verification stage.

Unlike approaches that make decisions solely on the basis of the number of
supported paths, our verification component analyzes the specific content
of the query to identify the type of claim being asked. The system
distinguishes standard causal questions, claims about direct edges, claims
about removing a mediator, claims about mediator necessity, and
counterfactual claims whose expected outcome is not explicitly specified.
This analysis uses only the query text and the retrieved structure; it does
not use \texttt{question\_type}, reference answers, or benchmark evaluation
labels.

To verify reasoning paths, we construct a Legal Structural Causal Model
(LegalSCM). LegalSCM is a deterministic reasoning model with the
three-valued state space
\(\{\mathrm{TRUE},\mathrm{FALSE},\mathrm{UNKNOWN}\}\). For each causal path,
the model checks whether the corresponding legal mechanisms can activate
the entire factual path. When a query requires counterfactual analysis,
LegalSCM performs a hard intervention \(do(M=\mathrm{FALSE})\) on mediator
\(M\), severs the mediator's incoming mechanisms, and re-infers the outcome.
The study also employs a simpler baseline in which the mediator is deleted
from the graph and the system checks whether an alternative path to the
outcome remains. This baseline is termed \emph{path ablation}; it captures
reachability after node deletion and is not interpreted as a full hard
intervention.

The verification result is passed to intent-aware answer generation. Rather
than assuming that every question asks for the final effect, the system
supports multiple output forms, including the final effect, initial cause,
bridge event, step-by-step explanation, yes/no answer, branching structure,
convergent structure, and claim verification. The extractive answer is
constructed from the primary causal path, prioritizes the complete effect
text of the final rule, and cites only the evidence actually used.

Experiments are conducted on a corpus of 2,884 rules and 1,814 events
normalized from the Vietnamese Penal Code. The resulting graph contains
5,134 nodes, 14,420 edges, and 112 two-hop causal chains. The system is
evaluated on a synthetic benchmark comprising 250 questions that cover
forward reasoning, reverse reasoning, intermediate-event identification,
chain explanation, yes/no questions, branching, and convergence. Across the
full benchmark, the system achieves a Rule Recall@5 of \(69.93\%\), an Event
Recall@5 of \(74.66\%\), a verification accuracy of \(92.00\%\), an Answer
Token F1 of \(68.69\%\), and a Citation F1 of \(62.86\%\). The gap between
Top-1 Exact Path (\(35.20\%\)) and Oracle Exact Path (\(71.20\%\)) shows that
many correct paths appear in the candidate set but are not ranked first.

The paired ablation between LegalSCM and path ablation reveals no difference
in the quality metrics on the current benchmark. The two methods attain the
same verification accuracy, Answer Token F1, and Citation F1. However,
LegalSCM requires an average of \(3.890\) seconds per sample, compared with
\(1.205\) seconds for path ablation, corresponding to a time cost
approximately \(3.23\) times higher. This result provides no evidence that
structural verification improves accuracy on the current benchmark;
instead, the observed benefit of LegalSCM lies in its clearer intervention
semantics and its ability to produce traces of the factual world,
intervention assignment, and counterfactual outcome.

The main contributions of this study are as follows:

\begin{itemize}
    \item We propose an end-to-end pipeline that combines vector memory, a
    legal causal graph, multi-hop retrieval, query-aware verification, and
    answer generation grounded in statutory provisions.

    \item We construct a three-valued LegalSCM to validate factual causal
    paths and perform hard interventions on intermediate events within the
    scope of normalized legal rules.

    \item We construct a synthetic evaluation set comprising 250 questions
    that cover multiple forms of forward, reverse, mediator-based,
    branching, and convergent legal reasoning.

    \item We conduct a paired ablation between structural SCM and
    node-deletion path ablation, thereby identifying a null quality result
    on the current benchmark and quantifying the time cost of structural
    verification.

    \item We analyze the remaining limitations, particularly causal-path
    ranking errors, low performance on reverse reasoning, and the lack of
    counterfactual samples with intervention ground truth.
\end{itemize}

The remainder of this paper is organized as follows.
Section~\ref{sec:related-work} presents related work on RAG, graph-based RAG,
causal reasoning, and legal question answering. Section~\ref{sec:problem}
formulates the problem and defines the scope of the claims.
Section~\ref{sec:method} describes the pipeline and LegalSCM.
Section~\ref{sec:data} presents the corpus and benchmark. The experimental
setup and results are presented in Sections~\ref{sec:experimental-setup}
and~\ref{sec:results}, respectively. Section~\ref{sec:discussion} analyzes
errors and discusses the results. Finally, the limitations and conclusion
are presented in Sections~\ref{sec:limitations} and~\ref{sec:conclusion}.
